\chapter{Einführung}
\label{chap:chapter-1}

In den letzten Jahren hat sich der Einsatz Künstlicher Intelligenz immer weiter verbreitet \cite{chatgpt}.
Zu den wichtigsten Methoden gehören Deep-Neural-Networks wie zum Beispiel Convolutional Neural Networks (CNNs) und Residual Networks (ResNets).
Diese werden unter anderem zur Bildklassifikation eingesetzt und sind mittlerweile in der Lage dazu, Ergebnisse mit sehr hoher Genauigkeit zu erzielen \cite{GIUDICI2023101560, he2016resnet, tan2021efficientnetv2}.

Ein großer Nachteil dieser Modelle ist allerdings die schlechte Nachvollziehbarkeit der Ergebnisse, wodurch sie auch als \enquote{Black-Box-Modelle} bezeichnet werden.
Explainable Artificial Intelligence (XAI) kann daher genutzt werden, um Nutzern KI-Entscheidungen, die auf komplexen Modellen basieren, einfacher darzustellen \cite{blackbox}.

Das Ziel dieser Projektarbeit war es, eine interaktive grafische Benutzeroberfläche zu entwickeln, mithilfe dieser die Bildklassifikation verständlicher gemacht werden kann. So soll das Tool \textit{Visual Explainable AI} Studierenden den Einstieg in das Fachgebiet \textit{Explainable Artificial Intelligence} erleichtern.

Es wurde zudem eine Fallstudie durchgeführt, die den Einfluss unseres Tools auf Studierende und deren Verständnis der Grenzen von KI-Modellen, sowie die Nachvollziehbarkeit und den Einfluss von XAI Erklärungen untersucht. \\

In dieser Hinsicht wurde die vorliegende Projektarbeit folgendermaßen strukturiert:
\begin{itemize}
    \item \textit{Kapitel 2: \nameref{chap:chapter-2}} erklärt die Grundlagen von neuronalen Netzwerken und XAI.
    \item \textit{Kapitel 3: \nameref{chap:chapter-3}} beschreibt die Anforderungsanalyse und begründet die getroffenen Entscheidungen bei der Auswahl der vortrainierten CNN-Modelle und der XAI-Verfahren.
    \item \textit{Kapitel 4: \nameref{chap:chapter-4}} beschreibt das entworferne Softwarekonzept sowie die geplante Systemarchitektur und begründet die getroffenen Entscheidungen.
    \item \textit{Kapitel 5: \nameref{chap:chapter-5}} bietet einen Überblick über die wichtigsten Entscheidungen bei der Implementierung des Tools.
    \item \textit{Kapitel 6: \nameref{chap:chapter-6}} zeigt wie die Fallstudie zur Evaluierung des entwickelten Tools aufgebaut und durchgeführt wurde, sowie die erzielten Ergebnisse.
    \item \textit{Kapitel 7: \nameref{chap:chapter-7}} evaluiert die Ergebnisse der Fallstudie und diskutiert die Wirksamkeit des Tools.
    \item \textit{Kapitel 8: \nameref{chap:chapter-8}} fasst die wichtigsten Erkenntnisse der Arbeit zusammen.
\end{itemize}

